\documentclass[10pt,a4paper]{amsart}
\usepackage[margin=19mm]{geometry}
\usepackage{amsmath,amssymb,mathtools}
\usepackage[scr=rsfs,cal=euler]{mathalfa}
\usepackage{xfrac}
\usepackage[unicode,hidelinks,bookmarks=false]{hyperref}
\newcommand{\ie}{i.e.\ }
\newcommand{\cf}{cf.\ }
\newcommand{\resp}{respectively\ }
\newcommand{\Subsection}{\subsection}
\providecommand{\nobreakdash}{\nobreak\hbox{-}}
\setlength{\emergencystretch}{2em}
\multlinegap0pt
\usepackage{bm}
\usepackage{bbm}
\usepackage{dsfont}
\usepackage{xspace}
\usepackage{cancel}
\usepackage{mathtools}
\usepackage{amsthm}
\makeatletter
\@ifundefined{theorem}{\newtheorem{theorem}{Theorem}}{}
\@ifundefined{definition}{\newtheorem{definition}[theorem]{Definition}}{}
\@ifundefined{proposition}{\newtheorem{proposition}[theorem]{Proposition}}{}
\makeatother
\def\normal{\mathcal{N}}
\title[On some log-concavity properties of the Alexander-Conway and]{On some log-concavity properties of the Alexander-Conway and Links-Gould invariants}
\author{Matthew Harper, Ben-Michael Kohli, Jiebo Song, Guillaume Tahar}
\date{Source: arXiv:2509.16868v1 (CC BY 4.0)}
\begin{document}
\maketitle

\noindent\emph{Source and licence.} On some log-concavity properties of the Alexander-Conway and Links-Gould invariants. Version arXiv:2509.16868v1 (CC BY 4.0), distributed under the Creative Commons Attribution 4.0 International licence, as recorded in the arXiv licence metadata for this exact version and shown on the abstract page https://arxiv.org/abs/2509.16868v1. Full rights notice: licenses/arxiv-2509.16868-CC-BY-4.0.txt.\par\smallskip

\noindent\textit{Editorial scope.} Counted unit: the Definition whose source label is \texttt{None} in the article (source0.tex lines 1111--1113, source label \texttt{None}), delivered together with the article's own complete original proof of it (source0.tex lines 1115--1150, one \texttt{proof} environment of 36 lines). No mathematics is written, repaired or re-derived: statement and proof are the article's own text line for line. Required context retained uncounted: algebraic lemma (lines 450--452). Every definition, display and label of those lines is retained unchanged. Omitted from this package and not redistributed: 1--449, 453--1110, 1114--1114, 1151--1984. Nothing inside a retained excerpt is omitted or altered: every retained body is delivered exactly as the article's lines are written, so no omission receipt of any kind accompanies this package. Citations: the retained excerpts carry no citation command at all, so no bibliography entry of the article is transcribed and no reference is redistributed. Reference adaptation: none was needed and none was made; every reference inside the delivered unit resolves inside it, so no unresolved reference or citation is produced. Printed slips kept literal: no printed word, symbol, hypothesis or number of the counted statement or of its proof was altered. Layout and class: the article's own class is \texttt{amsart} with packages including fullpage, amssymb, amsmath, mathtools, bbm, tikz, tikz-cd, color, dsfont, nicematrix, fancybox, hyperref, slashed, verbatim; the wrapper is the lane's pinned \texttt{amsart} editorial wrapper at 10pt on A4, which restates the theorem-like environments the retained text needs and adds the article's own macro definitions that the retained text needs (\texttt{\textbackslash normal}), with extra wrapper packages (bm, bbm, dsfont, xspace, cancel, mathtools). The delivered numbering is the unit's own. Assets: no retained span contains a figure environment, a graphics-inclusion command, a PDF-page inclusion, a table or a TikZ picture; no image file is copied into this package, the package declares no asset dependency and no image file is redistributed with it. Licence: version arXiv:2509.16868v1 of this article is distributed under the Creative Commons Attribution 4.0 International licence, as recorded in the arXiv licence metadata for this exact version and shown on the abstract page https://arxiv.org/abs/2509.16868v1. A full rights notice is delivered at licenses/arxiv-2509.16868-CC-BY-4.0.txt. Characters: every retained excerpt is pure ASCII, so the delivered engine, XeLaTeX with its default Latin Modern text font, reports no missing character.
\begin{proposition}\label{algebraic lemma}
For $P(t)$ a \emph{one-variable} polynomial with nonnegative real coefficients, its sequence of coefficients is log-concave if and only if $Q(t,z) := \mathrm{Homog}(P)$ is denormalized Lorentzian.
\end{proposition}

\begin{definition}
    A homogeneous polynomial $P$ is called \emph{denormalized Lorentzian} if $\normal(P)$ is Lorentzian.
\end{definition}

\begin{proof}[Proof of Proposition~\ref{algebraic lemma}]
We write 
$$
P(t) = \sum\limits_{k = 0}^{d} c_k\, t^k \,, \quad x_k \geq 0 \,, \quad x_d \neq 0 \,.
$$ 
To translate the Lorentzian condition in terms of coefficients of the polynomial, we can now compute the partial derivatives of $$Q(t,z) = \normal\circ\mathrm{Homog}(P) = \sum\limits_{k = 0}^{d} \frac{c_k\, t^k z^{d-k}}{k!(d-k)!}  \, .$$
Set $a,b \in \mathbb{Z}$ such that $0 \leq a \leq d-2$ and $a + b = d-2$. Then
\begin{align*}
 \frac{\partial^a Q}{\partial t^a} (t,z)   %& = \sum\limits_{k = 0}^{d} \frac{x_k}{k!(d-k)!} k (k-1) \ldots (k-a+1) \, t^{k-a} z^{d-k} \\
  & =  \sum\limits_{k = a}^{d} \frac{c_k}{(k-a)!(d-k)!} \, t^{k-a} z^{d-k}  
\end{align*}
and
\begin{align*}
 \frac{\partial^{a+b} Q}{\partial z^b \partial t^a} (t,z)   %& = \frac{\partial^{d-2} Q}{\partial z^b \partial t^a} (t,z) \\
  %& =  \sum\limits_{k = a}^{d-b} \frac{c_k}{(k-a)!(d-k-b)!} \, t^{k-a} z^{d-k-b} \\ 
  & =  \sum\limits_{k = a}^{a+2} \frac{c_k}{(k-a)!(a+2-k)!} \, t^{k-a} z^{a+2-k}
  = \frac{c_a}{2} z^2 + c_{a+1} t z + \frac{c_{a+2}}{2} t^2 \,.
\end{align*}
The Hessian matrix for $\frac{\partial^{a+b} Q}{\partial z^b \partial t^a} (t,z)$ relative to the derivatives $\partial z, \partial t$ is:
\[
\begin{pmatrix}
    c_a                          & c_{a+1} \\
 c_{a+1}                      & c_{a+2}
\end{pmatrix}
\]
with characteristic polynomial and discriminant
\begin{align*}
 \chi_a(\lambda)   %& = (c_a - \lambda)(c_{a+2}-\lambda)-c_{a+1}^2 \\
   =  \lambda^2 - (c_a + c_{a+2}) \lambda + c_a c_{a+2} - c_{a+1}^2 
   \qquad 
   \mbox{and}
   \qquad 
   \Delta_a = (c_a - c_{a+2})^2 + 4 c_{a+1}^2 \geq 0 \,.
\end{align*}
The Lorentzian property for $Q(t,z)$ means that the smaller root of $\chi_a$ is nonpositive. In other words, $c_a + c_{a+2} - \sqrt{\Delta_a} \leq 0$ for any $0 \leq a \leq d-2$. Since all the $c_i$ are nonnegative, this last inequality is equivalent to $c_a c_{a+2} \leq c_{a+1}^2$, which exactly means that the sequence $(c_i)$ of coefficients of $P$ is log-concave. This proves Proposition~\ref{algebraic lemma}.
\end{proof}

\end{document}
